\section{The majorant and the crossover at \texorpdfstring{$N/(\log N)^2$}{N/(log N)\^{}2}}
\label{app:boundary}

This appendix proves Lemma~\ref{lem:uniform-laguerre-bessel},
Theorem~\ref{thm:uniform-simplex-majorant} and the results behind
Remark~\ref{rem:crossover-scale}. The start is uniform.

\subsection{The majorant}

\begin{proof}[Proof of Lemma~\ref{lem:uniform-laguerre-bessel}]
By definition, $B_a(w)/w=\sum_{j=0}^{a-1}\binom{a-1}j w^j/(j+1)!$. So the
coefficient of $(aw)^j/[j!(j+1)!]$ in $B_a(w)/w$ is
$\binom{a-1}jj!/a^j=\prod_{h=1}^j(1-h/a)_+$, also for $j\ge a$, where it is $0$.
By Lemma~\ref{lem:clipped-product} its deficit from $1$ lies between $0$ and
$\sum_{h\le j}h/a=j(j+1)/(2a)$. Since $\sum_jj(j+1)x^j/[j!(j+1)!]=x\Phi(x)$
(Lemma~\ref{lem:bessel-law}(i)), we get
\[
0\le\Phi(aw)-\frac{B_a(w)}w
\le\frac1{2a}\sum_{j\ge0}\frac{j(j+1)(aw)^j}{j!(j+1)!}
=\frac w2\Phi(aw).\qedhere
\]
\end{proof}

\begin{proof}[Proof of Theorem~\ref{thm:uniform-simplex-majorant}]
\emph{Step 1. An integral form.} Let $0<u<1$ and $v=u/(1-u)$. Put
$M!=\int_0^\infty e^{-t}t^M\,dt$ into \eqref{eq:simplex-positive-sum} and sum
over each $m_i$ separately. All terms are nonnegative, so
\begin{equation}\label{eq:boundary-integral-form}
S_n(u)=\frac{(1-u)^{N-1}}{v}\int_0^\infty e^{-t}\prod_{i=1}^kB_{n_i}(vt)\,dt,
\end{equation}
with $B_a$ as in Lemma~\ref{lem:uniform-laguerre-bessel}.

\emph{Step 2. Comparison of the factors.} For $t>0$ put
$f_i(t)=B_{n_i}(vt)/(vt\,\Phi(n_ivt))\in[0,1]$. By
Lemma~\ref{lem:uniform-laguerre-bessel} we have $f_i(t)=(1-\theta_i)_+$ with
$0\le\theta_i\le vt/2$, so Lemma~\ref{lem:clipped-product} gives
$0\le1-\prod_if_i(t)\le kvt/2$. Write $\langle g\rangle$ for the mean of
$g(t)$ under the probability density on $(0,\infty)$ proportional to
$e^{-t}t^k\prod_i\Phi(n_ivt)$. Then \eqref{eq:boundary-integral-form} reads
$S_n(u)=\mathcal Q_n(u)\langle\prod_if_i\rangle$, and so
\[
0\le1-\frac{S_n(u)}{\mathcal Q_n(u)}=\Bigl\langle1-\prod_if_i\Bigr\rangle
\le\frac{kv}{2}\langle t\rangle.
\]
The bound $\Phi(x)\le e^{2\sqrt x}$ of Section~\ref{sec:multistate-model}
bounds the density by a polynomial times
$\exp(-t+2\sqrt{vt}\sum_i\sqrt{n_i})$, so all moments below exist.

\emph{Step 3. The mean of $t$.} We use $\psi(x)=x\Phi'(x)/\Phi(x)$ from
Lemma~\ref{lem:bessel-law}. Since $\frac d{dt}\log\Phi(n_ivt)=\psi(n_ivt)/t$,
\[
\frac{d}{dt}\Bigl[e^{-t}t^{k+1}\prod_i\Phi(n_ivt)\Bigr]
=e^{-t}t^k\prod_i\Phi(n_ivt)\Bigl(k+1-t+\sum_i\psi(n_ivt)\Bigr).
\]
The function in square brackets vanishes at $t=0$ and tends to $0$ as
$t\to\infty$, so integration gives
$\langle t\rangle=k+1+\sum_i\langle\psi(n_ivt)\rangle$. By
Lemma~\ref{lem:bessel-law}, $\psi(x)\le\sqrt x$. With
$\langle\sqrt t\rangle\le\sqrt{\langle t\rangle}$ and
$\Gamma=\sqrt v\sum_i\sqrt{n_i}$ this gives
$\langle t\rangle\le k+1+\Gamma\sqrt{\langle t\rangle}$. Since
$\Gamma\sqrt{\langle t\rangle}\le(\Gamma^2+\langle t\rangle)/2$, we get
$\langle t\rangle\le2(k+1)+v(\sum_i\sqrt{n_i})^2$. This proves the first
bound in \eqref{eq:uniform-simplex-majorant}, and Cauchy--Schwarz,
$(\sum_i\sqrt{n_i})^2\le kN$, gives the second.
\end{proof}

\subsection{The crossover}

Fix integers $k_1\ge2$ and $\ell\ge0$, and put $k=k_1+\ell$. We take $k_1$ large
counts $n_1,\ldots,n_{k_1}$ and $\ell$ rare counts $a_1,\ldots,a_\ell$, all
positive, and write $S_{(n,a)}$ for the ratio $S$ of this count vector. For
$\ell=0$ there are no rare counts, and the sums and products over $j$ below are
empty. Put
\[
N=\sum_in_i+\sum_ja_j,\quad M=\sum_in_i,\quad
L=\log N,\quad \beta_i=\frac{n_i}{M},\quad
\rho_j=\frac{a_jL^2}{N}.
\]
The $\beta_i$ and $\rho_j$ always refer to these actual counts. In this
appendix
\[
A=\sum_{i=1}^{k_1}\sqrt{\beta_i},\quad g=A^2-1,\quad
\eta=\frac{k_1-1}{2}+2\ell,\quad c_*=\frac{k-1}g,
\]
\[
D_{k_1}(\beta)=\frac{A^{k_1/2+1}(\prod_i\beta_i)^{-3/4}}
{(4\pi)^{(k_1-1)/2}},\qquad
\gamma_0=\frac{(2\eta+1)(2\eta+3)}{16A^2}-\frac{3}{16A}\sum_{i=1}^{k_1}\beta_i^{-1/2}.
\]
The constant $D_{k_1}(\beta)$ is $D(\beta)$ of
Theorem~\ref{thm:simplex-interior} on $k_1$ states, so it has the Kirchhoff
form of Proposition~\ref{prop:kirchhoff},
$D_{k_1}(\beta)=A^2(4\pi)^{-(k_1-1)/2}\sqrt{\tree(\beta)}/\prod_i\beta_i$. We use
$\psi(x)=x\Phi'(x)/\Phi(x)$ and its properties from
Lemma~\ref{lem:bessel-law}.

\begin{theorem}[a crossover between neighboring simplex strata]
\label{thm:simplex-boundary-crossover}
Fix $\varepsilon>0$ and $R<\infty$. Assume $\beta_i\ge\varepsilon$ and
$0<\rho_j\le R$. For $z>0$ put $x_j=A^2a_jz^2/N$ and
\[
\mathcal E(z)=-gz\sum_{j=1}^\ell\frac{a_j}{N}
+\frac{1}{A^2z}\Bigl(\eta\sum_{j=1}^\ell\psi(x_j)
+2\sum_{1\le i<j\le\ell}\psi(x_i)\psi(x_j)\Bigr)+\frac{\gamma_0}{z},
\]
so that $\mathcal E(z)=\gamma_0/z$ when $\ell=0$. Let $0<c_0<C_0<\infty$ be
fixed. Then, uniformly for such count vectors and $c_0L\le z\le C_0L$,
\begin{equation}\label{eq:simplex-boundary-profile}
S_{(n,a)}(z/N)
=D_{k_1}(\beta)A^{2\ell}\frac{z^\eta e^{gz}}{N^{k-1}}
\prod_{j=1}^\ell\Phi(x_j)\,e^{\mathcal E(z)}
\left(1+O\!\left(L^{-2}+\frac{L^2}{N}\right)\right).
\end{equation}
Also $|\mathcal E(z)|\le C_1(1+R)/L$ there, where $C_1$ depends only on
$\varepsilon$, $c_0$, $C_0$, $k_1$ and $\ell$. So
\eqref{eq:simplex-boundary-profile} also holds without the factor
$e^{\mathcal E(z)}$ if the error is replaced by $O(L^{-1}+L^2/N)$.

For large $N$, the right side of \eqref{eq:simplex-boundary-profile}
without its error factor equals $1$ at exactly one $\zeta_N$ in
$[(k-1)L/(2(k_1-1)),\,(k-1)L/\varepsilon]$. The unique positive solution $u_{n,a}$ of
$S_{(n,a)}(u)=1$ satisfies
\begin{equation}\label{eq:simplex-boundary-implicit-root}
Nu_{n,a}=\zeta_N+O\!\left(L^{-2}+\frac{L^2}{N}\right)
\end{equation}
and
\begin{align}
Nu_{n,a}
={}&c_*L-\frac\eta g\log L
-\frac1g\left\{
\log\bigl(D_{k_1}(\beta)A^{2\ell}\bigr)+\eta\log c_*
+\sum_{j=1}^\ell\log\Phi(A^2c_*^2\rho_j)\right\}
\notag\\
&\hspace{20mm}+O\!\left(\frac{\log L}{L}+\frac{L^2}{N}\right).
\label{eq:simplex-boundary-root}
\end{align}
The constants in the errors of \eqref{eq:simplex-boundary-profile} depend
only on $\varepsilon$, $R$, $k_1$, $\ell$, $c_0$ and $C_0$. Those in
\eqref{eq:simplex-boundary-implicit-root} and
\eqref{eq:simplex-boundary-root}, and the meaning of ``large $N$'', depend
only on $\varepsilon$, $R$, $k_1$ and $\ell$.
\end{theorem}

So each rare count multiplies the profile of the large counts by $A^2z^2/N$ and
by a Bessel factor $\Phi(x_j)$. At $u=z/N$ we have $T=z+O(L^2/N)$, the
logarithms of the right sides have bounded derivatives in $z$, and all errors
contain $L^2/N$. So the theorem also holds with $T$ in place of $z$ and of
$Nu_{n,a}$.

\begin{proof}
Let $u=z/N$ with $c_0L\le z\le C_0L$, and take $N$ large. Put
$v=u/(1-u)$, $Y=Mv$, $w_0=A\sqrt Y$, $q_0=\eta+3/2=2\ell+1+k_1/2$,
$s_a=\sum_ja_j/N$ and $\xi_j=a_jvw_0^2=A^2a_jMv^2$. Then $s_a\le\ell R/L^2$
and $Y=z(1-s_a)/(1-z/N)$, so $Y\asymp L$. Also $\xi_j\le X=4k_1C_0^2R$. When
$\ell=0$ we have $M=N$ and $s_a=0$, and the steps that involve the $\xi_j$ and
$X$ are empty. Note that $A^2=1+2\sum_{i<i'}\sqrt{\beta_i\beta_{i'}}$ lies in
$[1+2\varepsilon,k_1]$.

\emph{Step 1. The Bessel integral.} By
Theorem~\ref{thm:uniform-simplex-majorant}, with relative error
$O(v+Nv^2)=O(L^2/N)$, the ratio $S_{(n,a)}(u)$ equals $\mathcal Q_{(n,a)}(u)$.
Writing $\Phi(\beta_iYt)=I_1(2\sqrt{\beta_iYt})/\sqrt{\beta_iYt}$ for the large
counts, this is
\begin{equation}\label{eq:simplex-mixed-integral}
\frac{(1-u)^{N-1}v^{\ell-1}}{\prod_{i=1}^{k_1} n_i}
\int_0^\infty e^{-t}t^\ell
\prod_{i=1}^{k_1}\left[\sqrt{\beta_iYt}\,
I_1(2\sqrt{\beta_iYt})\right]
\prod_{j=1}^\ell\Phi(a_jvt)\,dt.
\end{equation}
Substitute $t=w^2$. The integrand becomes
\[
2w^{2\ell+1}e^{-w^2}\prod_{i=1}^{k_1}\Bigl[\sqrt{\beta_iY}\,w\,
I_1\bigl(2\sqrt{\beta_iY}\,w\bigr)\Bigr]\prod_{j=1}^\ell\Phi(a_jvw^2).
\]
Put $F(w)=w^{q_0}\prod_j\Phi(a_jvw^2)$, so that
$F(w_0)=A^{q_0}Y^{q_0/2}\prod_j\Phi(\xi_j)$.

\emph{Step 2. The main part.} On $I=\{|w-w_0|\le w_0/2\}$ the Bessel arguments
$2\sqrt{\beta_iY}\,w$ are at least $A\sqrt\varepsilon\,Y$. The expansion of
$I_1$ in Section~\ref{sec:multistate-model}, and
$-w^2+2A\sqrt Y\,w=A^2Y-(w-w_0)^2$, then show that on $I$ the integrand equals
\[
\frac{2(\prod_i\beta_i)^{1/4}Y^{k_1/4}}{(4\pi)^{k_1/2}}\,
e^{A^2Y-(w-w_0)^2}\,F(w)
\left(1-\frac{b_Y}w+O(Y^{-2})\right),\qquad
b_Y=\frac3{16\sqrt Y}\sum_i\beta_i^{-1/2}.
\]
We apply Lemma~\ref{lem:gaussian-phi-average} with $c_j=a_jv$, once with
$q=q_0$ to $F$ and once with $q=q_0-1$ to $F(w)/w$ for the term $b_Y/w$.
For the term $O(Y^{-2})$ we use its bound $F(w_0+y)=O(F(w_0))$. Since
$b_Y/w_0=O(1/Y)$, the integral over $I$ equals
\[
\frac{2\sqrt\pi(\prod_i\beta_i)^{1/4}Y^{k_1/4}}{(4\pi)^{k_1/2}}\,
e^{A^2Y}F(w_0)
\left(1+\frac{\mu_{q_0}}{4A^2Y}-\frac{3}{16AY}\sum_i\beta_i^{-1/2}
+O(Y^{-2})\right),
\]
with $\mu_{q_0}$ as in the lemma.

\emph{Step 3. The tails.} The part outside $I$ is negligible. The bounds
$I_1(x)\le e^x$ and $\Phi(x)\le e^{2\sqrt x}$ bound the integrand there by a
polynomial in $w$ and $Y$ times $\exp(A^2Y-(w-w_0)^2+2\delta w)$, where
$\delta=\sum_j\sqrt{a_jv}$, and $\delta w_0=\sum_j\sqrt{\xi_j}$ is
bounded. Completing the square gives $-(w-w_0-\delta)^2+2\delta w_0+\delta^2$,
and outside $I$ we have $|w-w_0-\delta|\ge3w_0/8$ for large $N$. Hence this
part is at most a power of $Y$ times $e^{A^2Y-9A^2Y/64}$, which is
$O(L^{-2})$ relative to the main term.

\emph{Step 4. Collecting the constants.} Since
$v^{\ell-1}/\prod_in_i=Y^{\ell-1}M^{-(k-1)}/\prod_i\beta_i$ and
$\ell-1+k_1/4+q_0/2=\eta$, the constants collect to
$D_{k_1}(\beta)A^{2\ell}Y^\eta M^{-(k-1)}$. Since $4q_0-6=4\eta$ and
$q_0(q_0-1)=(2\eta+1)(2\eta+3)/4$, the formula for $\mu_q$ in
Lemma~\ref{lem:gaussian-phi-average} gives
\[
\frac{\mu_{q_0}}{4A^2}-\frac{3}{16A}\sum_i\beta_i^{-1/2}
=\frac1{A^2}\sum_j\xi_j+\Theta(\xi),\quad
\Theta(\xi)=\frac1{A^2}\Bigl(\eta\sum_j\psi(\xi_j)
+2\sum_{i<j}\psi(\xi_i)\psi(\xi_j)\Bigr)+\gamma_0.
\]
Note that $\sum_j\xi_j/(A^2Y)=v\sum_ja_j$. Both $\sum_j\xi_j/A^2$ and
$\Theta(\xi)$ are bounded, so the bracket in Step 2 equals
$\exp(v\sum_ja_j+\Theta(\xi)/Y)(1+O(Y^{-2}))$. We have proved
\begin{equation}\label{eq:simplex-boundary-M-form}
S_{(n,a)}(u)=D_{k_1}(\beta)A^{2\ell}\frac{Y^\eta}{M^{k-1}}(1-u)^{N-1}
e^{v(A^2M+\sum_ja_j)}\prod_j\Phi(\xi_j)\,e^{\Theta(\xi)/Y}
\left(1+O\!\left(L^{-2}+\frac{L^2}{N}\right)\right).
\end{equation}

\emph{Step 5. From $(M,Y)$ to $(N,z)$.} Now use $u=z/N$. Since
$A^2M+\sum_ja_j=N(A^2-gs_a)$ and $vN=z+O(L^2/N)$, the exponent is
\[
v\Bigl(A^2M+\sum_ja_j\Bigr)+(N-1)\log(1-u)=gz-gzs_a+O(L^2/N).
\]
Also $Y^\eta M^{-(k-1)}=z^\eta N^{-(k-1)}(1+O(L^{-2}+L/N))$ and
$\xi_j=x_j(1-s_a)(1-z/N)^{-2}=x_j(1+O(L^{-2}+L/N))$. Since
$(\log\Phi)'(x)=\psi(x)/x\le1/2$, replacing $\xi_j$ by $x_j$ in $\Phi$
changes the product by a relative $O(L^{-2}+L/N)$. Since
$0\le\psi'\le1$, replacing $\Theta(\xi)/Y$ by $\Theta(x)/z$ changes it by
$O(L^{-3}+1/N)$. Finally $\mathcal E(z)=-gzs_a+\Theta(x)/z$, which proves
\eqref{eq:simplex-boundary-profile}.

\emph{Step 6. The bound on $\mathcal E$.} Note that $\psi(x_j)\le\sqrt{x_j}$,
$2\psi(x_i)\psi(x_j)\le\psi(x_i)^2+\psi(x_j)^2\le x_i+x_j$, and $A\ge1$. So
$|\mathcal E(z)|\le gzs_a+(\eta\sum_j\sqrt{x_j}+(\ell-1)\sum_jx_j+|\gamma_0|)/z$.
Here $gzs_a\le(k_1-1)C_0\ell R/L$, $x_j=A^2\rho_j(z/L)^2\le k_1C_0^2R$,
$z\ge c_0L$ and $\sqrt R\le1+R$, and $|\gamma_0|$ is bounded in terms of
$\varepsilon$, $k_1$ and $\ell$.

\emph{Step 7. The crossing.} Apply \eqref{eq:simplex-boundary-profile} with
$c_0=(k-1)/(2(k_1-1))$ and $C_0=(k-1)/\varepsilon$. Since
$g\in[2\varepsilon,k_1-1]$, we have $gc_0-(k-1)\le-(k-1)/2$ and
$gC_0-(k-1)\ge k-1$. Let $f(z)$ be the logarithm of
the right side of \eqref{eq:simplex-boundary-profile} without its error
factor. Then $f'(z)=g+(\eta+2\sum_j\psi(x_j))/z+\mathcal E'(z)$, since
$x_j'=2x_j/z$. Also $\mathcal E'(z)=O(L^{-2})$, because $gs_a=O(L^{-2})$,
$0\le\psi'\le1$ and the $x_j$ are bounded. So $f'\ge g/2\ge\varepsilon$
on $[c_0L,C_0L]$ for large $N$. Hence $f(c_0L)=(gc_0-(k-1))L+O(\log L)$ is
negative and $f(C_0L)=(gC_0-(k-1))L+O(\log L)$ is positive, and $f$ has exactly
one zero $\zeta_N$ in the interval. By \eqref{eq:simplex-boundary-profile} the
same signs hold for $\log S_{(n,a)}(z/N)$. The polynomial $S_{(n,a)}$ has
nonnegative coefficients, zero constant term and the positive coefficient
$k!$ at $u^{k-1}$, so the crossing is unique and lies in the interval. There
$|f(Nu_{n,a})|=O(L^{-2}+L^2/N)$, and $f'\ge\varepsilon$ gives
\eqref{eq:simplex-boundary-implicit-root}.

\emph{Step 8. The explicit expansion.} For \eqref{eq:simplex-boundary-root},
take $f(\zeta_N)=0$ and use $\mathcal E=O(L^{-1})$. This gives
\[
g\zeta_N+\eta\log\zeta_N=(k-1)L-\log(D_{k_1}(\beta)A^{2\ell})
-\sum_j\log\Phi\bigl(A^2\rho_j(\zeta_N/L)^2\bigr)+O(L^{-1}).
\]
First $\zeta_N=c_*L+O(\log L)$. The derivative of each
$\log\Phi(A^2\rho_j(z/L)^2)$ in $z$ is $2\psi/z=O(L^{-1})$, so the sum of
logarithms differs from its value at $\zeta_N=c_*L$ by $O(\log L/L)$.
Substituting this into $\log\zeta_N$ and using
\eqref{eq:simplex-boundary-implicit-root} proves
\eqref{eq:simplex-boundary-root}.
\end{proof}

\begin{remark}\label{rem:crossover-numbers}
The three terms of $\mathcal E$ have different sources. Replacing $M$ by $N$ in
the exponent $A^2Y$ gives a factor $\exp(-A^2z\sum_ja_j/N)$, and the Gaussian
average of the factors $\Phi$ gives back $\exp(z\sum_ja_j/N)$. Together they
give the first term. The second term comes from the same Gaussian average,
since the factors $\Phi$ and the power of $w$ shift the saddle point. The last
term does not depend on the sizes $a_j$, and for $\ell=0$ and equal $\beta_i$ it
is $-(k-1)/(8z)$, as in~\eqref{eq:simplex-uniform-ratio}.

Take two equal large counts and one rare count, so that $g=1$ and
$\eta=\frac52$, and let $z=\log N$. Then $x_1=2\rho$, and the first two terms of
$z\mathcal E$ are $-\rho+1.25\,\psi(2\rho)$, which nearly cancel at $\rho=1$. So
the first term alone is not the correction. The ratio of $S_{(n,a)}$ to the
right side of~\eqref{eq:simplex-boundary-profile} without $e^{\mathcal E}$ is
$1.04$ for $\rho=1$ and $N=2000$, and with the first term alone it would be
$1.19$. For $\rho=8$ it is $0.36$, $0.53$, $0.63$ and $0.70$ at $N=2\cdot10^3$,
$2\cdot10^4$, $2\cdot10^5$ and $2\cdot10^6$, and with the factor
$e^{\mathcal E}$ these become $0.52$, $0.70$, $0.79$ and $0.85$. For $\rho=1$
it becomes $0.90$. At $N=2000$ we have $L=7.6$, and the terms of order $L^{-2}$
and $L^2/N$ are comparable to $\gamma_0/L$, so this value does not improve. The
explicit
expansion~\eqref{eq:simplex-boundary-root} converges slowly. For $\rho=8$ and
$N=10^6$ it gives $6.58$, while $Nu_{n,a}=14.99$ and $\zeta_N=14.86$, since it
expands around $z=c_*L$ with $c_*=2$, while $Nu_{n,a}/L$ is only $1.08$
(numerically verified, see Section~\ref{sec:verify}).
\end{remark}

\begin{corollary}[invisibility of subcritical rare counts]
\label{cor:simplex-rare-invisibility}
Under the hypotheses of Theorem~\ref{thm:simplex-boundary-crossover}, if
every $a_j=o(N/L^2)$, the sum of the terms $\log\Phi$ in
\eqref{eq:simplex-boundary-root} is $o(1)$. So the sizes of the rare counts
change the expansion only by $o(1)$. Their number $\ell$ still changes its
leading and logarithmic coefficients.
\end{corollary}

\begin{proof}
Since $\Phi(0)=1$ and $0\le(\log\Phi)'\le1/2$ (Lemma~\ref{lem:bessel-law}),
we have $0\le\log\Phi(x)\le x/2$, and here $x=A^2c_*^2\rho_j\to0$.
\end{proof}

The scale $N/(\log N)^2$ needs at least two coordinates of order $N$. It does
not apply near a vertex, for example at an endpoint in the two-state case,
where only one coordinate is large and $g=0$
(Subsection~\ref{sec:vertex-completion}). The theorem gives no error bound
when some $\rho_j$ is unbounded.
